\section{Discussion}
\label{sec:discussion}

\paragraph{What a released ternary checkpoint reveals.}
Without training code, a released file still supports a surprising amount of
verification: the format, the rotation basis, the layout conventions, and the
code distribution are all measurable, and the public base makes the residual
measurable too. In this case the container turned out to be fully recoverable
from the artifact plus the base model, and the recovery later matched the
authors' published manifest exactly. Releasing format metadata is cheap and
disproportionately useful: it converts format reconstruction into verification
and lets independent parties spend their effort on the scientific question
instead.

\paragraph{Agreement is not quality.}
Ninety-two percent trit agreement sounds like a near-copy; it coexists with a
$1{,}270\times$ perplexity collapse. The mismatched codes are not noise around
a threshold, they are boundary placements that training put where the model
needs them, and they carry the model. Any pipeline that reports code agreement
as evidence of parity---ours included, before the payload-swap experiment---should treat the metric
as a container check and validate quality by a controlled swap or a held-out
evaluation.

\paragraph{Calibration hygiene is a first-class result.}
The public PTQ baseline reproduces its published figure on its calibration
passage and collapses by three to four orders of magnitude off it. This is not
a criticism of the method family but of the evaluation protocol: with
calibration-dependent methods, the calibration and evaluation sets must be
disjoint, and the disjointness must be stated. Our own protocol failures
(\S\ref{sec:pilot}) are the same disease in a different organ---an evaluation
set that was not actually held out---and they were caught only by an explicit
audit. We report the retraction because a forensic study's credibility rests on
its own error handling.

\paragraph{Fair accounting for ``about two bits per weight''.}
Effective bits per weight must count every parameter, including embeddings and
norms kept at higher precision. Ternary-only counts overstate compression: our
own otherwise ternary 1.7B recipe sat at an effective $4.53$\,bpw because the
embedding and head were not ternarized, while the released model ternarizes its
embedding, so a like-for-like comparison requires the effective number on both
sides. This matters for any claim that a PTQ method at ``2-bit'' matches a
trained ternary release.

\paragraph{Negative results, recorded.}
Several routes were eliminated: error-compensated quantization in the correct
basis, stochastic calibration resonance, continuous attractor potentials,
alternating projection, local per-layer distillation, entropy-based data
selection, and the mirror map. Each is recorded with its evidence and the cost
to revisit (Appendix~\ref{app:failures}). Two of them are informative in their
own right: the direction of the GPTQ results (compensation moves codes away)
and the mirror-map shape argument (a single-well force cannot populate a
two-well target) are general observations about ternary training, not just
statements about this model.

\paragraph{Open questions.}
Three remain. First, a full 27B end-to-end QAT/KD run, priced but not funded,
is the only route to a reproduction of the released numbers. Second, a
same-ruler comparison of trained ternary versus PTQ at matched effective bit
width, on the same base and corpus, is still missing from the public record.
Third, the extension of the forensics to mixture-of-experts models---routing
drift under ternarization and where to place corrections---is started in the
artifact but out of scope here; it is the natural next specimen, because
routing adds a failure mode that dense models do not have.
